% from ega2-1-fr.tex
\oldpage[II]{112}

En vertu de (\hyperref[I.5.5.12-fr]{I, 5.5.12}), il suffit de démontrer (i),
(ii) et (iii). Pour prouver qu'une immersion fermée $X\to S$ est finie, on
peut se borner au cas où $S=\operatorname{Spec}(A)$, et tout revient alors à
remarquer qu'un anneau quotient $A/\mathfrak{J}$ est un $A$-module monogène.
Pour prouver que le composé de deux morphismes finis (resp. entiers) $X\to Y$,
$Y\to Z$ est fini (resp. entier), on peut encore supposer $Z$ (et par suite
$X$, $Y$ (\hyperref[II.1.3.4-fr]{1.3.4})) affines, et alors l'assertion
équivaut à dire que si $B$ est une $A$-algèbre finie (resp. entière) et $C$
une $B$-algèbre finie (resp. entière), alors $C$ est une $A$-algèbre finie
(resp. entière), ce qui est immédiat. Enfin, pour démontrer (iii), on peut se
borner au cas où $S=Y$, puisque $X_{(S')}$ s'identifie à
$X\times_Y Y_{(S')}$ (\hyperref[I.3.3.11-fr]{I, 3.3.11})~; on peut en outre
supposer que $S=\operatorname{Spec}(A)$, $S'=\operatorname{Spec}(A')$~; alors
$X$ est affine d'anneau $B$ (\hyperref[II.1.3.4-fr]{1.3.4}), $X_{(S')}$
affine d'anneau $A'\otimes_A B$, et il suffit de remarquer que si $B$ est une
$A$-algèbre finie (resp. entière), $A'\otimes_A B$ est une $A'$-algèbre finie
(resp. entière).

Notons aussi que si $X$ et $Y$ sont deux $S$-préschémas qui sont finis (resp.
entiers) au-dessus de $S$, leur \emph{somme} $X\amalg Y$ est un préschéma fini
(resp. entier) au-dessus de $S$, car cela revient à voir que si $B$ et $C$
sont deux $A$-algèbres finies (resp. entières) sur $A$, il en est de même de
$B\times C$.

\begin{corollary}[6.1.6]
\label{II.6.1.6-fr}
Si $X$ est un préschéma entier (resp. fini) au-dessus de $S$, alors, pour tout
ouvert $U\subset S$, $f^{-1}(U)$ est entier (resp. fini) au-dessus de $U$.
\end{corollary}

C'est un cas particulier de (\hyperref[II.6.1.5-fr]{6.1.5, (iii)}).

\begin{corollary}[6.1.7]
\label{II.6.1.7-fr}
Soit $f:X\to Y$ un morphisme fini. Alors, pour tout $y\in Y$, la fibre
$f^{-1}(y)$ est un schéma algébrique fini sur $k(y)$, et a fortiori son
espace sous-jacent est discret et fini.
\end{corollary}

En effet, en tant que $k(y)$-préschéma, $f^{-1}(y)$ est identifié à
$X\times_Y\operatorname{Spec}(k(y))$ (\hyperref[I.3.6.1-fr]{I, 3.6.1}), qui
est fini au-dessus de $\operatorname{Spec}(k(y))$
(\hyperref[II.6.1.5-fr]{6.1.5, (iii)})~; c'est donc un schéma affine dont
l'anneau est une algèbre de rang fini sur $k(y)$
(\hyperref[II.6.1.4-fr]{6.1.4}). La proposition résulte alors de
(\hyperref[I.6.4.4-fr]{I, 6.4.4}).

\begin{corollary}[6.1.8]
\label{II.6.1.8-fr}
Soient $X$, $S$ deux préschémas intègres, $f:X\to S$ un morphisme dominant. Si
$f$ est entier (resp. fini), le corps $R(X)$ des fonctions rationnelles sur
$X$ est algébrique (resp. algébrique de degré fini) sur le corps $R(S)$ des
fonctions rationnelles sur $S$.
\end{corollary}

En effet, soit $s$ le point générique de $S$~; le $k(s)$-préschéma
$f^{-1}(s)$ est entier (resp. fini) au-dessus de
$\operatorname{Spec}(k(s))$ (\hyperref[II.6.1.5-fr]{6.1.5, (iii)}), et
contient par hypothèse le point générique $x$ de $X$~; l'anneau local de $x$
dans $f^{-1}(s)$ étant égal à $k(x)$ (\hyperref[I.3.6.5-fr]{I, 3.6.5}) est un
anneau local d'une algèbre entière (resp. finie) sur $k(s)$
(\hyperref[II.6.1.4-fr]{6.1.4}), d'où le corollaire.

\begin{remark}[6.1.9]
\label{II.6.1.9-fr}
L'hypothèse que $g$ est \emph{séparé} est essentielle pour la validité de
(\hyperref[II.6.1.5-fr]{6.1.5, (v)})~: en effet, si $Y$ n'est pas séparé sur
$Z$, l'identité $1_Y$ est le morphisme composé
$Y\xrightarrow{\Delta_Y}Y\times_ZY\xrightarrow{p_1}Y$, mais $\Delta_Y$ n'est
pas un morphisme entier, comme il résulte de
(\hyperref[II.6.1.10-fr]{6.1.10})~:
\end{remark}

\begin{proposition}[6.1.10]
\label{II.6.1.10-fr}
Tout morphisme entier est universellement fermé.
\end{proposition}

Soit $f:X\to Y$ un morphisme entier~; en vertu de
(\hyperref[II.6.1.5-fr]{6.1.5, (iii)}), il suffit de montrer que $f$ est
fermé. Soit $Z$ une partie fermée de $X$~; il existe un sous-préschéma de $X$
ayant

